\SourceLeaf{63}
per $\left(1-\dfrac{r}{\nu}\right)$, et multiplicatio per $\dfrac{r}{\nu}$ ita confici possit, ut
ille, secundum $n$ differentiatus, denuo in $n$ ducatur et per $\nu\,\mathrm{d}n$ dividatur,
relationem sequentem praebet: esse $B^{(r)}$ term. init.
\[
= A^{(r)}\text{ term. init.} - \frac{n\,{}^{*}\!\mathrm{d}A^{(r)}\text{ term. init.}}{\nu\,\mathrm{d}n}.
\]
Prorsus eadem est relatio, quae in terminis generalibus obtinet. Est nimirum ipsius
$A^{(r)}$ terminus $h^{\mathrm{tus}}$
\[
= \frac{1}{2^{r+2h-1}}\,{}^{r+2h}\mathfrak{B}^{(r+h)} . {}^{\nu}\mathfrak{B}^{(r+2h)}\,n^{r+2h}.
\]
Fit igitur, mutato $\nu$ in $\nu-1$, ipsius $B^{(r)}$ terminus $h^{\mathrm{tus}}$
\[
= \frac{1}{2^{r+2h-1}}\,{}^{r+2h}\mathfrak{B}^{(r+h)} . {}^{\nu-1}\mathfrak{B}^{(r+2h)}\,n^{r+2h},
\]
qui solo factore
\[
\frac{\nu-(r+2h)}{\nu} = 1 - \frac{r+2h}{\nu}
\]
ad priorem expressionem adiecto, ex ista sponte procedit. Cum autem idem sit, ipsius
$A^{(r)}$ terminum $h^{\mathrm{tum}}$ per $\dfrac{r+2h}{\nu}$ multiplicare, ac eundem, secundum
$n$ differentiatum, per $\nu\,\mathrm{d}n$ dividere et in $n$ denuo ducere, habebimus
\[
B^{(r)}\text{ term. }h = A^{(r)}\text{ term. }h - \frac{n\,{}^{*}\!\mathrm{d}A^{(r)}\text{ term. }h}{\nu\,\mathrm{d}n},
\]
quae relatio, cum pro omnibus sine discrimine membris obtineat, praebet aequationem
\[
B^{(r)} = A^{(r)} - \frac{n\,{}^{*}\!\mathrm{d}A^{(r)}}{\nu\,\mathrm{d}n}.
\]
\par\noindent\hspace*{\parindent}$\varepsilon$) Per hanc igitur scalam relationis $(1+n\cos\varphi)^{\nu-1}$ ad $(1+n\cos\varphi)^{\nu}$
ita reducitur, ut ex istius coefficientibus, qui in hanc expressionem cadant, differentiatione
facillime eliciantur. Licet adeo, cum ipsius $(1+n\cos\varphi)^{-1}$ coefficientes formulis
finitis circumscripti in potestate sint, ad exponentes negativos maiores sensim adscendere,
eoque omnino potestates negativas easque integras coefficientibus finitis exhibere.

Non tamen series infinitae, pro coefficientibus hisce ante inventae, inutiles sunt. Est iis
locus, simulac $\nu$ fuerit numerus fractus, quo casu, ex \textsc{Euleri} iudicio, in his
speculationibus gravissimo, formulae finitae nulla ratione sperari
